\documentclass[11pt]{article}

\usepackage{amsmath,amssymb,amsfonts}
\usepackage{hyperref}
\usepackage{geometry}
\usepackage{graphicx}
\usepackage{booktabs}
\usepackage{listings}
\usepackage{enumitem}

\geometry{margin=1in}

% --------------------
% Minimal macros
% --------------------
\newcommand{\dd}{\mathrm{d}}
\newcommand{\ii}{\mathrm{i}}
\newcommand{\ee}{\mathrm{e}}
\newcommand{\RR}{\mathbb{R}}
\newcommand{\CC}{\mathbb{C}}
\newcommand{\norm}[1]{\left\lVert #1 \right\rVert}
\newcommand{\abs}[1]{\left\lvert #1 \right\rvert}

\title{Actionable Plan: 1D Heat Equation Demo via Hybrid CV--DV LCHS\\
in \texttt{hybridlane}}
\author{}
\date{\today}

\begin{document}
\maketitle

\section{Goal, deliverables, and success criteria}

\subsection*{Goal}
Implement the 1D heat equation example using the hybrid CV--DV LCHS approach in \texttt{hybridlane}.

\subsection*{Primary deliverable}
A runnable example script (suggested path in the \texttt{hybridlane} repo):
\[
\texttt{examples/heat\_equation\_1d\_cvdv\_lchs.py}
\]
that:
\begin{enumerate}[itemsep=2pt,topsep=2pt]
\item builds a discretized heat-equation generator $A$ for $N=2^m$ grid points,
\item constructs the CV--DV LCHS circuit (state prep $\rightarrow$ hybrid evolution $\rightarrow$ postselection/projection),
\item produces (postselected) output vector $\tilde{\mathbf u}(t)$,
\item compares to the classical reference $\mathbf u(t)=\exp(-At)\mathbf u(0)$,
\item prints an error report and unit-test-friendly diagnostics.
\end{enumerate}

\subsection*{Definition of success}
Choose a small, deterministic reference case (default below: Dirichlet BCs and $f(x)=\sin(\pi x/L)$). For $m=2$ ($N=4$) or $m=3$ ($N=8$), demonstrate:
\[
\frac{\|\tilde{\mathbf u}(t)-\mathbf u(t)\|_2}{\|\mathbf u(t)\|_2} \le \epsilon_{\mathrm{target}}
\]
with a documented error budget. A good initial target is $\epsilon_{\mathrm{target}}=10^{-2}$ and then tighten.

\subsection*{Source alignment}
Primary references used for this numerical experiment:
\begin{itemize}[itemsep=2pt,topsep=2pt]
\item \texttt{res\_base/Hybrid\_CV\_DV\_LCHS.pdf} (problem mapping and CV--DV setup).
\item \texttt{res\_refs/Quantum algorithm for linear non-unitary dynamics with near-optimal dependence on all parameters.pdf} (kernel/LCHS theory).
\item \texttt{res\_refs/Constant-Factor Improvements in Quantum Algorithms for Linear Differential Equations.pdf} (kernel-shape guidance).
\item \texttt{res\_refs/2407.10381v3.pdf} and in-repo gate tests (hybrid gate conventions).
\end{itemize}

\subsection*{Current code status (important)}
The current implementation uses \textbf{Model B} in \texttt{heat1d\_lchs\_cv\_dv.py}
(\texttt{--model B} is default):
\[
\ket{\psi_{\mathrm{post}}^{(B)}}=\sum_n C_n\ket{\psi_n},\quad
p_{\mathrm{succ}}^{(B)}=\norm{\psi_{\mathrm{post}}^{(B)}}_2^2,\quad
\rho_B=\frac{\ket{\psi_{\mathrm{post}}^{(B)}}\bra{\psi_{\mathrm{post}}^{(B)}}}{p_{\mathrm{succ}}^{(B)}}.
\]
where \(\ket{\psi_n}\) is the subnormalized DV branch from Fock component \(n\), with
\(p_n=\norm{\psi_n}_2^2\).

Legacy Model A material is retained in the appendix for comparison only.

\section{Mathematical specification}

\subsection{Continuous PDE model and analytic solution}

\paragraph{PDE.}
Let $u(x,t)$ denote temperature on a rod of length $L>0$. Let $\alpha>0$ be thermal diffusivity.
The 1D heat equation (paper Sec.~3.5.2, Eq.~(41)--(43)) is
\begin{equation}
\frac{\partial u(x,t)}{\partial t} = \alpha \frac{\partial^2 u(x,t)}{\partial x^2}, \quad 0<x<L,\; t\ge 0.
\end{equation}

\paragraph{Default boundary condition (Dirichlet).}
\begin{equation}
u(0,t)=0,\qquad u(L,t)=0,\qquad t\ge 0.
\end{equation}

\paragraph{Initial condition.}
\begin{equation}
u(x,0)=f(x),\qquad 0<x<L.
\end{equation}

\paragraph{Closed-form solution (Dirichlet).}
Define Fourier sine coefficients
\begin{equation}
b_n := \frac{2}{L}\int_0^L f(x)\sin\!\left(\frac{n\pi x}{L}\right)\dd x,\qquad n=1,2,\dots
\end{equation}
Then the analytic solution is
\begin{equation}
u(x,t) = \sum_{n=1}^{\infty} b_n \sin\!\left(\frac{n\pi x}{L}\right)\exp\!\left(-\alpha \left(\frac{n\pi}{L}\right)^2 t\right).
\end{equation}

\paragraph{Default reference case (recommended for unit tests).}
Set $f(x)=\sin(\pi x/L)$, which yields $b_1=1$ and $b_{n\ne 1}=0$, hence:
\begin{equation}
u(x,t)=\sin\!\left(\frac{\pi x}{L}\right)\exp\!\left(-\alpha \left(\frac{\pi}{L}\right)^2 t\right).
\end{equation}
This gives a clean, single-mode reference.

\subsection{Spatial discretization to a linear ODE}

\paragraph{Grid.}
Choose \(N=2^m\) interior points to amplitude-encode into \(m\) qubits.
For Dirichlet BCs, set grid spacing
\begin{equation}
h := \frac{L}{N+1},\qquad x_j := jh,\quad j=1,\dots,N.
\end{equation}
Paper baseline (Sec.~3.5.2) uses a 6-point stencil with 4 interior points (\(N=4,m=2\)).

\paragraph{State vector.}
Define $\mathbf u(t)\in\RR^N$ by $\mathbf u_j(t)\approx u(x_j,t)$.

\paragraph{Discrete Laplacian (Dirichlet).}
Define \(T\in\RR^{N\times N}\) as the tridiagonal matrix:
\begin{equation}
T := \begin{bmatrix}
2 & -1 \\
-1 & 2 & -1 \\
& \ddots & \ddots & \ddots \\
&& -1 & 2 & -1 \\
&&& -1 & 2
\end{bmatrix}.
\end{equation}
Then define the generator
\begin{equation}
A := \frac{\alpha}{h^2} T.
\end{equation}

\paragraph{Resulting ODE.}
\begin{equation}
\frac{\dd}{\dd t}\mathbf u(t) = -A\mathbf u(t),\qquad \mathbf u(0)=\mathbf u_0.
\end{equation}
For \(N=4\) this matches Eq.~(45) in the working paper with \(A=\alpha T/h^2\).

\paragraph{Theoretical discrete solution.}
\begin{equation}
\mathbf u(t)=\exp(-At)\,\mathbf u_0.
\end{equation}
\paragraph{Pauli decomposition for the 4-point case.}
In paper notation (Eq.~(46)):
\begin{equation}
T=2I_0I_1-I_0X_1-\frac12(X_0X_1+Y_0Y_1),
\end{equation}
equivalently,
\begin{equation}
T=2(I\otimes I)-(I\otimes X)-\frac12(X\otimes X+Y\otimes Y).
\end{equation}

\subsection{CV--DV LCHS target expression}

\paragraph{Cartesian decomposition (general form).}
Given $A=L+\ii H$, define
\begin{equation}
L:=\frac{A+A^\dagger}{2},\qquad H:=\frac{A-A^\dagger}{2\ii}.
\end{equation}
For diffusion-like dynamics we assume $L\succeq 0$ (stable).

\paragraph{LCHS continuous representation.}
Using the working-paper notation, choose
\begin{equation}
g(k):=\frac{f(k)}{1-ik},
\end{equation}
and
\begin{equation}
\exp\!\left(-\int_0^t A(s)\dd s\right)
=\int_{\RR} g(k)\left[\mathcal{T}\exp\!\left(-\ii\int_0^t (kL(s)+H(s))\dd s\right)\right]\dd k.
\end{equation}
For time-independent $A$, this simplifies to
\begin{equation}
\exp(-At)=\int_{\RR} g(k)\,\exp\!\left(-\ii t(kL+H)\right)\dd k.
\end{equation}

\paragraph{Improved kernel family (used in your working paper).}
Pick $\beta\in(0,1)$ and define
\begin{equation}
f(k)=\frac{\ee^{2\beta}}{2\pi \ee\, (1+\ii k)^\beta},\qquad
g(k)=\frac{f(k)}{1-\ii k}.
\end{equation}

\paragraph{CV--DV ``sandwich'' theorem target.}
Let $\hat{x}$ be the qumode position quadrature operator (quantum optics convention).
We want a circuit implementing the hybrid unitary
\begin{equation}
U(t):=\exp\!\left(-\ii t(\hat{x}\otimes L + I_{\mathrm{osc}}\otimes H)\right).
\end{equation}
Choose two normalized qumode states \(|\psi\rangle_{\mathrm{osc}}\) (prepared) and \(|\varphi\rangle_{\mathrm{osc}}\) (postselected)
with position wavefunctions \(\psi(x)\) and \(\varphi(x)\) such that
\begin{equation}
\varphi^*(x)\psi(x)=g(x).
\end{equation}
Then the induced (unnormalized) DV state is
\begin{equation}
|\mathbf u(t)\rangle_{\mathrm{DV}}
=
\left(\langle \varphi|_{\mathrm{osc}}\otimes I\right)\,U(t)\,\left(|\psi\rangle_{\mathrm{osc}}\otimes|\mathbf u_0\rangle_{\mathrm{DV}}\right)
=
\left[\int_{\RR} g(x)\,\exp\!\left(-\ii t(xL+H)\right)\dd x\right]|\mathbf u_0\rangle_{\mathrm{DV}}.
\end{equation}
In the ideal setting, this equals $\exp(-At)|\mathbf u_0\rangle$.

\paragraph{Heat equation specialization.}
For the heat equation, $A$ is real symmetric so $H=0$ and $L=A$.
Thus the target reduces to
\begin{equation}
|\mathbf u(t)\rangle
=
\left[\int_{\RR} g(x)\,\exp(-\ii t x A)\dd x\right]|\mathbf u_0\rangle.
\end{equation}

\section{Implementation architecture: classical vs quantum responsibilities}

\subsection{Classical preprocessing responsibilities}
Given $(\alpha,L,t,N=2^m)$ and initial function $f(x)$:
\begin{enumerate}[itemsep=3pt,topsep=3pt]
\item Build \(A=(\alpha/h^2)T\in\RR^{N\times N}\).
\item Build a reference solution $\mathbf u_{\mathrm{ref}}(t)=\exp(-At)\mathbf u_0$ (SciPy).
\item Obtain a Pauli decomposition on $m$ qubits:
\[
A = \sum_{j=1}^{J} c_j P_j,
\]
where each $P_j$ is a Pauli string on $m$ qubits and $c_j\in\RR$.
For small $m$, use automated helper utilities; for larger $m$, exploit operator structure.
\item Choose LCHS kernel hyperparameter $\beta\in(0,1)$ and define $g(\cdot)$.
\item Choose the two squeezing parameters from the paper:
\begin{itemize}[itemsep=2pt,topsep=2pt]
\item \(r\): squeezing for postselection state \(|\varphi_r\rangle\).
\item \(r'\): squeezing used in the prepared-branch basis.
\end{itemize}
\item Define \(\sigma=e^r,\ \sigma'=e^{r'}\), and
\[
\gamma=\frac{1}{\sigma'^2}-\frac{1}{\sigma^2}=e^{-2r'}-e^{-2r},\qquad r'<r.
\]
Compute prepared-state coefficients \(\{C_n\}_{n=0}^{N_{\max}-1}\) in the paper Eq.~(12) form:
\begin{equation}
C_n
\propto
\sqrt{\frac{\sigma}{\sigma'}}\,
\frac{1}{\sqrt{2^n n!}}
\int_{\RR} H_n(x/\sigma')\,g_\beta(x)\,e^{-\gamma x^2}\,\dd x.
\end{equation}
In code this integral is evaluated by Gauss--Hermite quadrature with the same normalization convention as
\texttt{lchs\_coefficients}. The improved-kernel branch is
\begin{equation}
g_\beta(x)=\frac{\exp\!\big(-(1+ix)^\beta\big)}{C_\beta(1-ix)},
\qquad
C_\beta=2\pi e^{-2^\beta},\ \beta\in(0,1).
\end{equation}
\item Normalize by \(\ell_2\): \(C_n\leftarrow C_n/\norm{C}_2\). The runtime mixture weights are \(p_n=|C_n|^2\).
\end{enumerate}

\subsection{Quantum circuit responsibilities}
With one qumode wire (call it \texttt{"m"}) and $m$ qubits (wires \texttt{0..m-1}):
\begin{enumerate}[itemsep=3pt,topsep=3pt]
\item Prepare the DV initial state $|\mathbf u_0\rangle$ (for this repository, a computational-basis state from \texttt{init\_q0, init\_q1}).
\item Prepare the qumode state $|\psi\rangle_{\mathrm{osc}}$ (details in Sec.~\ref{sec:stateprep}).
\item Apply the hybrid evolution $U(t)=\exp(-\ii t\,\hat{x}\otimes A)$ by Trotterization over Pauli terms (Sec.~\ref{sec:trotter}).
\item Postselect/projection: project the qumode onto $|\varphi\rangle_{\mathrm{osc}}$ (Sec.~\ref{sec:projection}).
\item Return the (postselected) DV state vector and the success probability.
\end{enumerate}

\section{Operator ledger: math definitions and hybridlane mapping}

\subsection{Primitive CV operators and conventions}
We assume the quantum-optics convention:
\[
\hat{x}=\frac{a+a^\dagger}{2},\quad \hat{p}=\frac{a-a^\dagger}{2\ii},\quad a=\hat{x}+\ii\hat{p}.
\]
This matches the Gaussian wavefunction convention used in the CV LCHS derivations.

\subsection{Hybridlane operator definitions (must match the math)}
Use these \texttt{hybridlane} operators:
\begin{itemize}[itemsep=2pt,topsep=2pt]
\item \texttt{hqml.Displacement(a, phi, wires="m")} implements
$\exp(\alpha a^\dagger-\alpha^* a)$ with $\alpha=ae^{\ii\phi}$.
\item \texttt{hqml.Squeezing(r, phi, wires="m")} implements
$S(\zeta)=\exp\left[\frac12(\zeta^*a^2-\zeta(a^\dagger)^2)\right]$ with $\zeta=re^{\ii\phi}$.
\item \texttt{hqml.ConditionalDisplacement(a, phi, wires=[q,\,"m"])} implements
$CD(\alpha)=\exp\left[\sigma_z(\alpha a^\dagger-\alpha^* a)\right]$ with $\alpha=ae^{\ii\phi}$.
\end{itemize}

\subsection{Implementation note}
The numerical experiment is defined by the \texttt{hybridlane} gate semantics used in this document and in the tests.
Treat low-level simulator details as opaque unless a gate-level unit test fails.

\section{State preparation and projection}
\label{sec:stateprep}

\subsection{Postselection state $|\varphi_r\rangle$}
Use a squeezed vacuum for projection:
\[
|\varphi_r\rangle := S(r,0)|0\rangle.
\]
Its position-space wavefunction is
\[
\varphi_r(x)=\left(\frac{2}{e^{2r}\pi}\right)^{1/4}\exp\!\left(-\frac{x^2}{e^{2r}}\right)
=\left(\frac{2}{\sigma^2\pi}\right)^{1/4}\exp\!\left(-\frac{x^2}{\sigma^2}\right).
\]
In the implemented circuit this is done by applying \(S(-r,0)\) to the qumode and then measuring the vacuum projector.
Equivalently:
\[
\langle \varphi_r|\psi_{\mathrm{final}}\rangle
=
\langle 0|S^\dagger(r,0)|\psi_{\mathrm{final}}\rangle,
\]
which is implemented by inverse squeezing followed by vacuum projection.

\subsection{Prepared state in this repository}
From \(\varphi_r^*(x)\psi_r(x)=g(x)\), the formal preparation wavefunction is
\[
\psi_r(x)=\mathcal N\,g(x)\exp\!\left(\frac{x^2}{\sigma^2}\right),\qquad \sigma=e^r.
\]
Given \(\varphi_r\) is Gaussian, \(\psi_r\) is generally non-Gaussian (especially for the near-optimal kernel with \(\beta\in(0,1)\)).

The implementation uses a truncated squeezed-Fock expansion:
\[
|\psi_r\rangle\approx\sum_{n=0}^{N_{\max}-1} C_n\,S(r',0)|n\rangle.
\]
The coefficient vector \(C_n(r,r',\beta)\) is used by the current Model-B path.

\subsection{Implemented CV preparation path (actual runtime)}
\paragraph{Mode B (default): coherent branch aggregation.}
For each active level \(n\), the circuit prepares \(S(r',0)\ket{n}\), runs the same hybrid evolution,
and extracts the subnormalized postselected DV branch \(\ket{\psi_n}\).
Then:
\[
\ket{\psi_{\mathrm{post}}^{(B)}}=\sum_n C_n\ket{\psi_n},\qquad
p_{\mathrm{succ}}^{(B)}=\norm{\psi_{\mathrm{post}}^{(B)}}_2^2.
\]
This preserves cross-term interference in \(n\)-space through the complex \(C_n\) phases.

\paragraph{Physical-circuit note.}
Current Model B is a coherent \emph{numerical} reconstruction from branch amplitudes.
It matches the coherent target map in simulation, but it is not yet a single-shot
hardware circuit that directly injects \(\sum_n C_n\ket{n}\).

\section{Hybrid Hamiltonian evolution by Trotterization}
\label{sec:trotter}

\subsection{Target unitary}
For heat equation, target:
\[
U(t)=\exp(-\ii t\,\hat{x}\otimes A),\qquad A=\sum_{j=1}^{J} c_j P_j.
\]
Thus:
\[
U(t)=\exp\left(-\ii t\sum_{j=1}^{J} c_j(\hat{x}\otimes P_j)\right).
\]

\subsection{First-order Trotter step}
Let $n_{\mathrm{Trot}}\in\mathbb{N}$ and $\delta t:=t/n_{\mathrm{Trot}}$.
Use:
\[
U(t)\approx \left(\prod_{j=1}^{J}\exp(-\ii \delta t\, c_j\,\hat{x}\otimes P_j)\right)^{n_{\mathrm{Trot}}}.
\]

\subsection{Implementing $\exp(-\ii \lambda\,\hat{x}\otimes P)$ using controlled displacement}
Key synthesis idea:
\[
CD(\alpha)=\exp\left[\sigma_z(\alpha a^\dagger-\alpha^*a)\right]
=\exp\left(+\ii(2\,\mathrm{Im}\,\alpha)\,\sigma_z\hat{x}\right)\cdot\exp\left(-\ii(2\,\mathrm{Re}\,\alpha)\,\sigma_z\hat{p}\right)
\times e^{i\varphi_{\mathrm{BCH}}},
\]
where \(e^{i\varphi_{\mathrm{BCH}}}\) is a global phase from non-commutation of \(\hat{x},\hat{p}\).
Choosing $\alpha = -\ii\lambda/2$ yields:
\[
CD(-\ii\lambda/2)=\exp(-\ii\lambda\,\sigma_z\hat{x}).
\]

To implement $\exp(-\ii \lambda\,\hat{x}\otimes P)$ for a multi-qubit Pauli string $P$:
\begin{enumerate}[itemsep=2pt,topsep=2pt]
\item Conjugate $P$ into a single-qubit $Z$ via Clifford basis changes and a parity-compute circuit.
\item Apply \texttt{hqml.ConditionalDisplacement(a, phi)} with $\alpha=-\ii\lambda/2$ on the parity qubit and qumode.
\item Uncompute the parity and undo basis changes.
\end{enumerate}

\section{Projection, success probability, and classical readout}
\label{sec:projection}

\subsection{Postselection}
In the current implementation we use projector-conditioned Pauli expectations, not direct full-statevector readout.
Let $\Pi_\phi$ be the CV projector used for postselection (implemented as inverse preparation + vacuum projector).
Measure:
\[
p_{\mathrm{succ}} = \ev{\Pi_\phi},
\qquad
m_{ab}=\ev{\Pi_\phi\,(P_a\otimes P_b)}.
\]
Reconstruct the conditional DV density matrix:
\[
\rho_{\mathrm{post}}=\frac{1}{4\,p_{\mathrm{succ}}}\sum_{a,b\in\{I,X,Y,Z\}} m_{ab}\,(P_a\otimes P_b).
\]
For vector-style PDE comparison, use the principal-eigenvector proxy of \(\rho_{\mathrm{post}}\).

\subsection{Convert back to a PDE comparison vector}
In this repository the DV register is compared to the finite-difference target via
\(\rho_{\mathrm{post}}\) and its principal eigenvector proxy; the PDE-priority metric is reported on
\(\sqrt{p_{\mathrm{succ}}}\,\psi_{\mathrm{principal}}\).

\section{Parameter-to-metric model for \texorpdfstring{$(r,r',\beta)$}{(r,r',beta)}}
\label{sec:param_model}

This section uses the paper parameters \((r,r',\beta)\).

\subsection{Coefficient map}
Define
\[
\gamma(r,r')=e^{-2r'}-e^{-2r},\qquad \gamma>0\ \Leftrightarrow\ r'<r.
\]
For fixed truncation \(N_{\max}\),
\[
C_n(r,r',\beta)\propto
\int_{\RR}
\exp\!\big[-\gamma(r,r')x^2\big]\,
g_\beta(x)\,
\frac{H_n(x/e^{r'})}{\sqrt{\sqrt{\pi}\,e^{r'}\,2^n n!}}\dd x,
\]
then normalize:
\[
\sum_{n=0}^{N_{\max}-1}\abs{C_n}^2=1,\qquad
w_n(r,r',\beta):=\abs{C_n(r,r',\beta)}^2.
\]
In the current runtime model, \((r,r',\beta)\) influence downstream metrics through both \(C_n\) and the
component responses \((\ket{\psi_n},p_n)\), because \(r,r'\) are also used directly in the component circuit
(pre-squeezing and postselection squeezing).

\subsection{Success probability}
Let \(p_n(r,r')\) be per-Fock success probabilities from the component circuit at fixed
phase and Trotter depth.
Then the coherent Model-B success probability is
\begin{equation}
p_{\mathrm{succ}}^{(B)}(r,r',\beta)
=
\left\|
\sum_{n=0}^{N_{\max}-1} C_n(r,r',\beta)\,\ket{\psi_n(r,r')}
\right\|_2^2.
\label{eq:psucc_rrb}
\end{equation}

\subsection{Coherent success structure (Model B)}
For each \(n\), let the subnormalized DV branch output be \(\ket{v_n}\) with
\(\norm{v_n}_2^2=p_n\), and define \(\ket{\chi_n}:=\ket{v_n}/\sqrt{p_n}\).
\[
\ket{u_B}=\sum_n C_n\ket{v_n}
=\sum_n C_n\sqrt{p_n}\ket{\chi_n},
\]
\[
p_{\mathrm{succ}}^{(B)}=\norm{u_B}_2^2
=\sum_n |C_n|^2p_n
\;+\!\!\sum_{n\neq m} C_n C_m^*\sqrt{p_np_m}\,\langle\chi_m|\chi_n\rangle.
\]
This decomposition makes explicit the diagonal and interference contributions.
Legacy Mode-A contrast is provided in the appendix.

\subsection{Algorithmic error function \texorpdfstring{$\epsilon_{\mathrm{alg,err}}(r,r',\beta)$}{epsilon_alg_err(r,r',beta)}}
\paragraph{Ideal coherent definition (paper-level).}
Let
\[
F_{r,r',\beta}(A,t)
:=
\int_{\RR} G_{r,r',\beta}(x)\,\exp(-\ii xtA)\,\dd x,
\]
where \(G_{r,r',\beta}\) is the effective spectral kernel induced by preparation/postselection.
Define
\begin{equation}
\epsilon_{\mathrm{alg,err}}^{\mathrm{ideal}}(r,r',\beta)
:=
\frac{\norm{\big(F_{r,r',\beta}(A,t)-e^{-At}\big)\mathbf u_0}_2}{\norm{e^{-At}\mathbf u_0}_2}.
\label{eq:eps_alg_ideal}
\end{equation}

\paragraph{Numerical algorithmic definition used in this repository (Model B).}
\[
\ket{u_B(r,r',\beta)}=\sum_n C_n(r,r',\beta)\ket{v_n(r,r')},\qquad
p_{\mathrm{succ}}^{(B)}=\norm{u_B}_2^2,
\]
\[
\rho_{\mathrm{alg}}^{(B)}(r,r',\beta)
=
\frac{\ket{u_B(r,r',\beta)}\!\bra{u_B(r,r',\beta)}}{p_{\mathrm{succ}}^{(B)}(r,r',\beta)}.
\]
Let \(\psi_{\max}(\rho)\) be the principal eigenvector of \(\rho\). Then:
\begin{equation}
\epsilon_{\mathrm{alg,err}}^{\mathrm{num},B}(r,r',\beta)
:=
\frac{\norm{
\mathbf u_{\mathrm{ref}}
-\sqrt{p_{\mathrm{succ}}^{(B)}(r,r',\beta)}\,
\psi_{\max}\!\big(\rho_{\mathrm{alg}}^{(B)}(r,r',\beta)\big)
}_2}{\norm{\mathbf u_{\mathrm{ref}}}_2},
\quad
\mathbf u_{\mathrm{ref}}=e^{-At}\mathbf u_0.
\label{eq:eps_alg_num}
\end{equation}
Equation~\eqref{eq:eps_alg_num} is the concrete function aligned with the script's PDE-priority proxy.

\subsection{Numerical implementation PDE solution error}
Using measured moments and reconstructed \(\hat\rho_{\mathrm{post}}\),
\begin{equation}
\epsilon_{\mathrm{PDE}}^{\mathrm{impl}}(r,r',\beta)
:=
\frac{\norm{
\mathbf u_{\mathrm{ref}}
-\sqrt{\hat p_{\mathrm{succ}}(r,r',\beta)}\,
\psi_{\max}\!\big(\hat\rho_{\mathrm{post}}(r,r',\beta)\big)
}_2}{\norm{\mathbf u_{\mathrm{ref}}}_2}.
\label{eq:eps_impl}
\end{equation}

Practical decomposition:
\[
\epsilon_{\mathrm{PDE}}^{\mathrm{impl}}
\lesssim
\epsilon_{\mathrm{alg,err}}^{\mathrm{ideal}}
+\epsilon_{\mathrm{fock}}
+\epsilon_{\mathrm{trot}}
+\epsilon_{\mathrm{tomography}}
+\epsilon_{\mathrm{num}}.
\]
For parameter sweeps, \((r,r',\beta)\) primarily move
\(\epsilon_{\mathrm{alg,err}}^{\mathrm{num}}\), \(p_{\mathrm{succ}}\), and \(\epsilon_{\mathrm{fock}}\) through
both coefficient weights \(w_n\) and component responses \((p_n,\rho_n)\).

\subsection{Directional effects of \texorpdfstring{$r,r',\beta$}{r,r',beta}}
\begin{itemize}[itemsep=3pt,topsep=3pt]
\item \textbf{\(r,r'\) via \(\gamma\):} larger \(r-r'\) means larger \(\gamma\), narrower \(k\)-envelope, usually lower effective mode count \(n_{\mathrm{eff}}\), less truncation stress, but possible spectral bias increase in \(\epsilon_{\mathrm{alg,err}}^{\mathrm{ideal}}\).
\item \textbf{\(r'\) basis scaling:} \(r'\) also rescales Hermite argument \(k/e^{r'}\), shifting Fock-weight placement; this changes \(p_{\mathrm{succ}}\) and error non-monotonically.
\item \textbf{\(\beta\) kernel shape:} \(\beta\) changes complex spectral weighting \(g_\beta(k)\), controlling how different eigenvalues of \(A\) are matched. Coupling to \(\gamma\) is strong: if \(\gamma\) is large (narrow \(k\)-window), \(\beta\)-sensitivity weakens.
\end{itemize}

\subsection{Asymptotic arguments}
With \(\theta=(r,r',\beta)\), \(\gamma(\theta)=e^{-2r'}-e^{-2r}>0\), and
\[
C_n(\theta)\propto\int_{\RR} e^{-\gamma(\theta)x^2} g_\beta(x)\,\Phi_n(x;e^{r'})\,\dd x,
\quad w_n(\theta)=|C_n(\theta)|^2,
\]
the current model metrics are
\[
p_{\mathrm{succ}}(\theta)=\sum_n w_n(\theta)\,p_n(r,r'),
\]
\[
\epsilon_{\mathrm{alg,err}}^{\mathrm{num}}(\theta)
=\frac{\norm{u_{\mathrm{ref}}-\sqrt{p_{\mathrm{succ}}(\theta)}\,
\psi_{\max}(\rho_{\mathrm{alg}}(\theta))}_2}{\norm{u_{\mathrm{ref}}}_2}.
\]

\paragraph{Regime A: \(\gamma\to\infty\) (large \(r-r'\)).}
The Gaussian factor localizes the \(k\)-integral to width \(O(\gamma^{-1/2})\) around \(k=0\).
Consequences: lower \(n_{\mathrm{eff}}\), smaller truncation tail, typically larger \(p_{\mathrm{succ}}\), but potentially higher spectral bias
(\(\epsilon_{\mathrm{alg,err}}^{\mathrm{ideal}}\) may increase if the required spectral support is broader).

\paragraph{Regime B: \(\gamma\to 0^+\) (\(r'\to r^{-}\)).}
The envelope broadens in \(k\), exciting more Fock modes.
Consequences: \(n_{\mathrm{eff}}\) increases, tail mass grows, truncation pressure rises, and \(p_{\mathrm{succ}}\) often drops due to weight spreading.

\paragraph{Regime C: kernel tail \(|k|\to\infty\).}
For \(\beta\in(0,1)\),
\[
|g_\beta(k)|\sim \frac{\exp\!\left(-|k|^\beta\cos(\pi\beta/2)\right)}{|k|}.
\]
The \(\beta\)-dependence is non-monotone in practical finite-\(k\) windows (because both \(|k|^\beta\) and
\(\cos(\pi\beta/2)\) vary with \(\beta\)). Asymptotic coherent-kernel guidance (often favoring \(\beta\approx 0.7\!-\!0.9\))
does not automatically transfer to Mode-A objectives, where dephasing and success-probability scaling are explicit.

\paragraph{End-to-end scaling form.}
A practical asymptotic decomposition is
\[
\epsilon_{\mathrm{PDE}}^{\mathrm{impl}}
\lesssim
\epsilon_{\mathrm{alg,err}}^{\mathrm{ideal}}(\theta)
+O(n_{\mathrm{steps}}^{-1})
+O\!\left(\sqrt{\tau_{N_{\max}}(\theta)}\right)
+\epsilon_{\mathrm{tomography}}
+\epsilon_{\mathrm{num}},
\]
where \(\tau_{N_{\max}}(\theta)=\sum_{n\ge N_{\max}}|C_n(\theta)|^2\).
This gives the optimization logic: first minimize \(\epsilon_{\mathrm{alg,err}}^{\mathrm{num}}\), then enforce \(p_{\mathrm{succ}}\) and tail constraints.

\subsection{PDE-priority optimization objective}
A practical single-objective form is
\begin{equation}
J(r,r',\beta)
=
\epsilon_{\mathrm{alg,err}}^{\mathrm{num}}(r,r',\beta)
+\lambda_p\frac{[p_{\min}-p_{\mathrm{succ}}(r,r',\beta)]_+}{p_{\min}}
+\lambda_t\frac{[\mathrm{tail}(r,r',\beta)-\tau_{\max}]_+}{\tau_{\max}},
\label{eq:objective_rrb}
\end{equation}
which encodes the priority order: PDE error first, success probability second.

\subsection{Hard feasibility constraints}
For theorem-constrained search, evaluate only points satisfying:
\begin{align}
&r' < r-\Delta_{\min}, \\
&\gamma(r,r')=e^{-2r'}-e^{-2r}\ge \gamma_{\min}, \\
&0<\beta<1, \\
&\tau_{N_{\max}}(r,r',\beta)\le \tau_{\max}, \\
&n_{\mathrm{eff}}(r,r',\beta)=\frac{1}{\sum_n |C_n(r,r',\beta)|^4}\le n_{\mathrm{eff},\max}.
\end{align}
These constraints prevent invalid squeeze geometry and control truncation complexity before objective ranking.

\section{Error budget and parameter selection guide}

Define a target end-to-end error $\epsilon_{\mathrm{target}}$ and allocate:
\[
\epsilon_{\mathrm{target}}
\ge
\epsilon_{\mathrm{space}}
+\epsilon_{\mathrm{kernel}}
+\epsilon_{\mathrm{prep}}
+\epsilon_{\mathrm{proj}}
+\epsilon_{\mathrm{fock}}
+\epsilon_{\mathrm{trot}}.
\]

\subsection{Recommended staged procedure}
\begin{enumerate}[itemsep=3pt,topsep=3pt]
\item \textbf{Stage 0 (gate sanity):} validate Displacement, ConditionalDisplacement, and Squeezing definitions by matrix tests at small cutoff.
\item \textbf{Stage 1 (LCHS identity scalar test):} set $A=\lambda I$ on 1 qubit, verify the sandwich yields $e^{-\lambda t}$.
\item \textbf{Stage 2 (PDE small-N):} run $m=2$ ($N=4$), compare to $\exp(-At)\mathbf u_0$.
\item \textbf{Stage 3 (near-optimal kernel):} switch to $\beta\in(0,1)$ and prepared non-Gaussian $|\psi_r\rangle$.
\end{enumerate}

\subsection{How to pick $\beta$, $r$, $r'$, and cutoff}
\begin{itemize}[itemsep=3pt,topsep=3pt]
\item \textbf{If implementing coherent Mode B (paper target):} start near Cauchy baseline \(\beta\approx 0.95\),
then explore \(\beta\in[0.7,0.9]\) as suggested by near-optimal-kernel numerics.
\item For historical Model-A-specific tuning behavior, see Appendix~\ref{sec:results}.
\item Increase $r$ to make the postselection state more sharply peaked in $x$ (but track qumode energy and cutoff needs).
\item Tune $r'<r$ to reduce coefficient spread in the truncated expansion of $|\psi_r\rangle$.
\item Increase \texttt{max\_fock\_level} until projection/cutoff errors stop dominating (diagnose via tail probability beyond cutoff).
\end{itemize}

\section{Unit test blueprint}

Create tests under:
\[
\texttt{tests/test\_gate\_math.py}
\]

\subsection*{Test list}
\begin{enumerate}[itemsep=3pt,topsep=3pt]
\item \textbf{Gate-definition tests}
\begin{enumerate}[itemsep=2pt,topsep=2pt]
\item Displacement matrix equality between hybridlane and analytic $D(\alpha)$.
\item ConditionalDisplacement generates $\exp(-\ii\lambda \sigma_z \hat{x})$ for purely imaginary $\alpha=-\ii\lambda/2$.
\item Squeezing-definition check: compare implemented squeezing matrix against analytic $S(\zeta)$; if mismatch, pin down the parameter transform and assert it.
\end{enumerate}
\item \textbf{Scalar LCHS test} ($A=\lambda I$): output scalar matches $e^{-\lambda t}$ within tolerance.
\item \textbf{End-to-end PDE test} ($m=2$): compare postselected output to $\exp(-At)\mathbf u_0$.
\item \textbf{Trotter convergence test}: error decreases as $n_{\mathrm{Trot}}$ increases (log-log fit optional).
\item \textbf{Success probability regression}: compare measured $p_{\mathrm{succ}}$ against classical $\|\mathbf u(t)\|_2^2$ in the high-accuracy limit.
\end{enumerate}

\section{Reproducible command protocol}
\paragraph{Step 0: gate-math sanity (fast fail).}
\begin{verbatim}
pytest -q tests/test_gate_math.py
\end{verbatim}

\paragraph{Step 1: phase calibration and baseline export.}
\begin{verbatim}
python heat1d_lchs_cv_dv.py --calibrate-phase \
  --output-json results/baseline_phase.json
\end{verbatim}
Use the selected \texttt{disp\_phase} from this step in all later runs.

\paragraph{Step 2: theorem-constrained sensitivity with fixed phase.}
\begin{verbatim}
python heat1d_lchs_sensitivity.py \
  --output-dir results/sensitivity \
  --n-r-target 7 --n-r-prime 7 --n-beta 9 \
  --disp-phase <PHASE_FROM_STEP1>
\end{verbatim}

\paragraph{Step 3: Trotter convergence at fixed \((r,r',\beta,\phi)\).}
\begin{verbatim}
for N in 10 25 50 100 200 500; do
  echo "=== n_steps=$N ==="
  python heat1d_lchs_cv_dv.py \
    --n-steps $N \
    --disp-phase <PHASE_FROM_STEP1> \
    --r-target <R> --r-prime <RP> --kernel-beta <KB> \
    --output-json results/trotter_n${N}.json
done
\end{verbatim}

\appendix
\section{Model-A Legacy Baseline}
\label{sec:results}

All results below use Mode~A (incoherent Fock-component aggregation),
$N=4$ interior points ($m=2$ qubits), Dirichlet BCs,
initial state $|\mathbf u_0\rangle=|01\rangle$, total time $t=1$,
Fock truncation $N_{\max}=32$, and $\hbar=2$.
These tables are retained as historical baseline data; current default runtime in code is Model~B.

\subsection{Phase calibration (Step~1)}
Phase calibration selected $\phi=0$ (x-dominant displacement),
with $\Delta x=+2$, $\Delta p=0$ at unit amplitude.
Candidate phases $\phi=\pm\pi/2$ gave PDE error $\approx 0.66$ vs $0.22$ at baseline,
confirming the correct quadrature branch.

\subsection{Sensitivity analysis (Step~2)}

A deterministic coarse grid ($7\times 7\times 9$ in $r,r',\beta$) followed by
local refinement around the top seeds identified the optimal parameters:
\begin{center}
\begin{tabular}{lcc}
\toprule
Parameter & Baseline & Optimized \\
\midrule
$r$ (target squeeze)    & 1.20 & 0.85 \\
$r'$ (prep squeeze)     & 0.30 & 0.03 \\
$\beta$ (kernel shape)  & 0.80 & 0.35 \\
$\phi$ (disp.\ phase)   & 0.0  & 0.0  \\
\midrule
$\gamma = e^{-2r'}-e^{-2r}$ & 0.458 & 0.759 \\
$n_{\mathrm{eff}}$           & 1.30  & 1.41  \\
tail mass                    & $9.5\times 10^{-6}$ & $1.3\times 10^{-6}$ \\
\bottomrule
\end{tabular}
\end{center}

\begin{center}
\begin{tabular}{lccr}
\toprule
Metric & Baseline & Optimized & Change \\
\midrule
PDE error        & 0.216 & 0.122 & $-44\%$ \\
Fidelity         & 0.948 & 0.957 & $+1\%$  \\
Post-selection probability & 14.2\% & 18.5\% & $+30\%$ \\
Pauli RMSE       & 0.138 & 0.085 & $-39\%$ \\
\bottomrule
\end{tabular}
\end{center}

Key observations from the grid:
\begin{itemize}[itemsep=2pt,topsep=2pt]
\item PDE error improves with moderate $r\approx 0.85$; larger $r$ degrades performance.
\item Minimal ancilla squeeze ($r'=0.03$) is consistently optimal across the grid.
\item Low $\beta=0.35$ outperforms higher $\beta$, consistent with the Mode-A analysis
      in Sec.~7.5: dephased aggregation couples $\epsilon_{\mathrm{PDE}}$ to
      $p_{\mathrm{succ}}^{(A)}$, not only directional fidelity.
\end{itemize}

\subsection{Trotter convergence (Step~3)}

At the optimized parameters $(r=0.85,\,r'=0.03,\,\beta=0.35,\,\phi=0)$:
\begin{center}
\begin{tabular}{rcccc}
\toprule
$n_{\mathrm{steps}}$ & PDE error & Fidelity & $p_{\mathrm{succ}}$ & Trotter bound \\
\midrule
 10  & 0.1067 & 0.9588 & 0.1822 & 0.050 \\
 25  & 0.1168 & 0.9577 & 0.1842 & 0.020 \\
 50  & 0.1203 & 0.9572 & 0.1848 & 0.010 \\
100  & 0.1220 & 0.9570 & 0.1851 & 0.005 \\
200  & 0.1229 & 0.9568 & 0.1852 & 0.0025 \\
500  & 0.1235 & 0.9568 & 0.1853 & 0.001 \\
\bottomrule
\end{tabular}
\end{center}

\paragraph{Interpretation.}
PDE error \emph{increases} slightly with $n_{\mathrm{steps}}$ and plateaus at $\approx 0.124$.
This is the opposite of Trotter-dominated behavior.
Applying the interpretation rule from the execution protocol:
\begin{quote}
\textbf{The bottleneck is non-Trotter.}
The $\approx 12\%$ PDE error floor is structural, arising from the
Mode-A incoherent aggregation mismatch ($\epsilon_{\mathrm{alg,err}}^{\mathrm{num}}$).
\end{quote}
The slight improvement at $n_{\mathrm{steps}}=10$ is attributable to partial cancellation
between Trotter error and algorithmic error (they subtract rather than add at low step count).
Fidelity and $p_{\mathrm{succ}}$ are essentially converged by $n_{\mathrm{steps}}=50$.

\subsection{Error budget summary}

At the converged operating point ($n_{\mathrm{steps}}=500$, optimized parameters):
\[
\epsilon_{\mathrm{PDE}}^{\mathrm{impl}} = 0.124
\approx
\underbrace{\epsilon_{\mathrm{alg,err}}^{\mathrm{num}}}_{\approx 0.12}
+ \underbrace{\epsilon_{\mathrm{trot}}}_{< 0.002}
+ \underbrace{\epsilon_{\mathrm{fock}}}_{< 0.001}
+ \underbrace{\epsilon_{\mathrm{num}}}_{< 0.001}.
\]
The dominant contribution is the Mode-A algorithmic error.
Fidelity $\mathcal{F}=0.957$ and purity $\mathrm{Tr}(\rho^2)=0.943$
indicate the postselected state direction is accurate but the
incoherent mixture introduces a $\sim 6\%$ mixed-state penalty.

\subsection{Implications for further improvement}

\begin{enumerate}[itemsep=3pt,topsep=3pt]
\item \textbf{Model B is now implemented as the default numerical path} and should be used for
      all new sweeps. The $\sim 12\%$ floor reported here is a Mode-A legacy artifact.
\item \textbf{Next target} is a single-shot physical circuit for coherent qumode injection
      $\sum_n C_n|n\rangle$ (rather than branch reconstruction), to close the remaining
      simulation-vs-hardware gap.
\item \textbf{If $\mathcal{F}\ge 0.95$ suffices as the success criterion},
      the current Mode-A implementation already meets it at the optimized parameters.
\item \textbf{Parameter tuning within Mode A is near-exhausted}:
      the sensitivity grid covered $7\times 7\times 9 = 441$ points plus refinement,
      and the Trotter sweep confirms the floor is not step-count limited.
\end{enumerate}

\section{Why incoherent aggregation is the bottleneck (legacy Mode-A analysis)}
\label{sec:coherence_gap}

The Trotter convergence data (Sec.~\ref{sec:results}) shows the $\approx 12\%$ PDE error
floor is entirely due to Mode~A incoherent aggregation.
This section explains the fundamental mechanism and proposes concrete replacements.

\subsection{What Mode~A breaks}

LCHS is fundamentally a \emph{coherent} linear combination:
\[
e^{-At} = \int g(x)\,e^{-itxA}\,\dd x.
\]
Mode~A replaces the coherent superposition $\sum_n C_n |v_n\rangle$ with a classical mixture
$\rho_A = \sum_n |C_n|^2 |\chi_n\rangle\!\langle\chi_n|$, discarding all cross-term interference
\[
\Delta p_{\mathrm{int}} = \sum_{n\neq m} C_n C_m^* \sqrt{p_n p_m}\,\langle\chi_m|\chi_n\rangle.
\]
This is not a tunable approximation error---it is a structural channel mismatch
that no parameter optimization within Mode~A can eliminate.

Two specific consequences:
\begin{enumerate}[itemsep=2pt,topsep=2pt]
\item \textbf{Lost interference.}
The kernel $g(x)$ is complex-valued (especially for the near-optimal $\beta\in(0,1)$ family).
Replacing coherent amplitude combination by classical averaging destroys the phase relationships
that make the integral converge to $e^{-At}$.
\item \textbf{Mixed-state output.}
Even when the directional fidelity $\mathcal{F}=0.957$ is high, the purity
$\mathrm{Tr}(\rho^2)=0.943<1$ means the output is a statistical mixture, not a pure state.
The PDE error metric penalizes both the direction error \emph{and} the norm mismatch
from $p_{\mathrm{succ}}^{(A)} \neq p_{\mathrm{succ}}^{(B)}$.
\end{enumerate}

\subsection{Coherent replacement options}
\label{sec:coherent_options}

\paragraph{Option A (recommended MVP): discrete coherent LCU with qubit index register.}

Discretize the LCHS integral into a finite sum and implement it as a standard LCU block-encoding:
\[
\sum_{j=-M}^{M} a_j\, U(t; k_j), \qquad U(t;k_j) = e^{-it k_j A}.
\]
Circuit skeleton:
\begin{enumerate}[itemsep=2pt,topsep=2pt]
\item \textbf{Prep}: $|0\rangle \mapsto \sum_j \sqrt{|a_j|/\|a\|_1}\,|j\rangle$
      (store $\arg(a_j)$ as phase gates).
\item \textbf{Select}: $\sum_j |j\rangle\!\langle j| \otimes U(t;k_j)$.
\item \textbf{Unprep} and postselect ancilla on $|0\rangle$.
\end{enumerate}
This keeps full coherence and gives controllable success probability.
For the heat equation with $H=0$, each $U(t;k_j) = e^{-itk_j A}$ uses the same
Trotter decomposition already implemented, parameterized by $k_j$.

\textbf{Why this is the right MVP}: it validates the entire ODE/PDE pipeline
with coherent linear combination before dealing with CV state-preparation subtleties.
Once correct, the index register can be upgraded from DV$\to$CV (Option~B).

\paragraph{Option B (paper target): continuous LCU via ancilla qumode + postselection.}

This is the architecture described in Sec.~\ref{sec:stateprep}--\ref{sec:projection},
executed coherently:
\begin{enumerate}[itemsep=2pt,topsep=2pt]
\item Prepare $|\psi\rangle_{\mathrm{osc}} \otimes |u_0\rangle_{\mathrm{DV}}$ as a \emph{single coherent state}.
\item Apply the joint unitary $e^{-it(\hat{x}\otimes A)}$.
\item Postselect $\langle\varphi|_{\mathrm{osc}}$, yielding
      $\big[\int \varphi^*(x)\psi(x)\,e^{-itxA}\,\dd x\big]|u_0\rangle = e^{-At}|u_0\rangle$.
\end{enumerate}
The key requirement is coherent injection of the truncated state
$|\psi\rangle \approx \sum_n C_n\, S(r',0)|n\rangle$ as a single quantum operation,
not as separate classical runs per Fock level.

\paragraph{Option C (fallback, not recommended): unbiased importance sampling.}

If coherent preparation is unavailable, the least-bad incoherent approach is importance sampling:
sample $k \sim p(k) \propto |g(k)|$ and reweight observables by $g(k)/p(k)$.
This can estimate \emph{specific observables} without bias, but does not produce the evolved state
in amplitudes and suffers large variance when $g$ has heavy tails or fast phase oscillation.
Not recommended for the hybridlane demonstration.

\subsection{Recommended implementation path}

\begin{center}
\begin{tabular}{clp{7cm}}
\toprule
Phase & Target & Deliverable \\
\midrule
1 & Option A (DV LCU) & Discrete coherent LCU with $2M+1$ quadrature points.
    Validate $e^{-At}|u_0\rangle$ end-to-end for the $N=4$ heat equation.
    Uses only qubit gates (no CV). \\
2 & Option B (CV LCU) & Replace the DV index register with a single qumode.
    Implement coherent $|\psi\rangle = \sum_n C_n S(r',0)|n\rangle$ injection.
    Reuse the same $U(t;k)$ backend from Phase~1. \\
\bottomrule
\end{tabular}
\end{center}

Both phases share the same core primitive: the parameterized Hamiltonian simulation
$U(t;k) = e^{-itkA}$, which is already implemented and Trotter-validated.
The transition from Phase~1 to Phase~2 only changes the ancilla encoding
(qubit register $\to$ qumode), not the system-register evolution.

\subsection{Minimum hybridlane primitives required}

For \textbf{Phase 1 (DV LCU)}:
\begin{itemize}[itemsep=2pt,topsep=2pt]
\item State preparation on ancilla register: arbitrary amplitude encoding of $\{a_j\}$.
\item Controlled-$U(t;k_j)$: select oracle conditioned on index register.
\item Ancilla postselection on $|0\rangle$.
\end{itemize}

For \textbf{Phase 2 (CV LCU)}, additionally:
\begin{itemize}[itemsep=2pt,topsep=2pt]
\item Coherent qumode state injection: $|0\rangle \mapsto \sum_n C_n |n\rangle$ (e.g., via
      SNAP gates, optimal control pulses, or sequential number-selective rotations).
\item Squeezing gate $S(r',0)$ applied after Fock superposition.
\item Joint CV--DV evolution $e^{-it\hat{x}\otimes A}$ (already implemented).
\item Inverse-squeeze + vacuum projection for postselection (already implemented).
\end{itemize}

\subsection{Relation to the existing codebase}

The current codebase uses Model~B by default and keeps Model~A as a baseline.
Mode~A remains useful as:
\begin{enumerate}[itemsep=2pt,topsep=2pt]
\item A \textbf{baseline} for validating Phase~1/Phase~2 outputs (the per-component
      $\{p_n, \rho_n\}$ data is correct; only the aggregation is incoherent).
\item A \textbf{diagnostic tool}: comparing Mode~A vs coherent LCU outputs isolates
      the interference contribution $\Delta p_{\mathrm{int}}$ experimentally.
\item A \textbf{lower bound}: Mode~A fidelity $\mathcal{F}=0.957$ should be exceeded
      by coherent implementations; if not, there is a circuit-level bug.
\end{enumerate}

\end{document}
